\section{Chain-aware anatomical assignment}
\label{sec:chain-aware}

Several downstream analyses in this report (the per-part correspondence
audit, the multi-level evaluation suite in \S\ref{sec:eval-suite})
depend on a prior step: deciding which anatomical structure a given surface
point belongs to. A naïve nearest-segment rule, which assigns each point to
whichever skeleton bone is geometrically closest, fails systematically on an
articulated body plan: a point on a middle leg can sit closer in raw
distance to the wrong leg, or to the body, than to its own leg, because the
rule treats every bone independently rather than as part of a connected
kinematic chain. A \textbf{chain-aware} rule, which resolves ambiguous
points by their position along the connected chain rather than by raw
distance, was tested against it on the same labelled template and adopted:
macro-averaged F1 improves from $0.656$ to $0.831$, concentrated exactly
where the naïve rule was weakest (middle-leg recall $0.349\to0.768$,
body/core recall $0.366\to0.663$). The lesson generalises beyond this one
rule: an articulated body plan is not a set of independent parts, and
geometric rules that treat it as one will fail systematically at the joints
between parts.

\begin{algbox}{2: Chain-aware anatomical assignment}
\begin{enumerate}[nosep,leftmargin=1.4em]
\item For each surface point $p$, compute distance to every bone segment.
\item Group segments into kinematic chains (each leg, antenna, mandible, body axis).
\item Assign $p$ to the chain whose segments are collectively closest, using
the chain's connected extent rather than the single nearest segment.
\item Within that chain, assign to the nearest segment.
\end{enumerate}
\end{algbox}
